
\documentclass{article}
\usepackage{colortbl}
\usepackage{makecell}
\usepackage{multirow}
\usepackage{supertabular}

\begin{document}

\newcounter{utterance}

\centering \large Interaction Transcript for game `clean\_up', experiment `1\_medium\_5obj\_ar', episode 1 with mistral{-}large{-}3{-}azure.
\vspace{24pt}

{ \footnotesize  \setcounter{utterance}{1}
\setlength{\tabcolsep}{0pt}
\begin{supertabular}{c@{$\;$}|p{.15\linewidth}@{}p{.15\linewidth}p{.15\linewidth}p{.15\linewidth}p{.15\linewidth}p{.15\linewidth}}
    \# & $\;$A & \multicolumn{4}{c}{Game Master} & $\;\:$B\\
    \hline

    \theutterance \stepcounter{utterance}  
    & & \multicolumn{4}{p{0.6\linewidth}}{
        \cellcolor[rgb]{0.9,0.9,0.9}{
            \makecell[{{p{\linewidth}}}]{
                \texttt{\tiny{[P1$\langle$GM]}}
                \texttt{أنا مدير لعبتك، وأنت تلعب لعبة تعاونية مع الشبكة التالية كلوحة لعب:} \\
\\ 
\texttt{```} \\
\texttt{ 1234567} \\
\texttt{╔══╤══╤═╗} \\
\texttt{║◌◌│T◌│◌║ 1} \\
\texttt{╟──┼──┴─╢ 2} \\
\texttt{║WI│◌◌◌◌║ 3} \\
\texttt{╟──┘◌◌◌◌║ 4} \\
\texttt{║◌◌◌H◌◌C║ 5} \\
\texttt{║◌◌┌─┬──╢ 6} \\
\texttt{║◌◌│◌│◌◌║ 7} \\
\texttt{╚══╧═╧══╝} \\
\\ 
\texttt{```} \\
\\ 
\texttt{* تعرض الحافة العلوية إحداثيات x التي تتزايد نحو اليمين، والحافة اليمنى إحداثيات y التي تتزايد نحو الأسفل.} \\
\texttt{* الكائنات التالية موضوعة عشوائياً على شبكتك: 'W', 'I', 'T', 'C', 'H'.} \\
\\ 
\\ 
\texttt{يرى اللاعب الآخر نسخة مختلفة من لوحة اللعب، حيث توضع الكائنات في مواقع عشوائية مختلفة. لا يمكنك رؤية لوحة اللاعب الآخر، ولا يمكنه رؤية لوحتك.} \\
\\ 
\\ 
\texttt{**الهدف:**} \\
\\ 
\texttt{يجب على كلا اللاعبين تحريك الكائنات على خلفية كل منهما بحيث تنتهي الكائنات المتطابقة عند الإحداثيات نفسها. عليك التواصل مع اللاعب الآخر للاتفاق على إعداد هدف مشترك.} \\
\\ 
\\ 
\texttt{**القواعد:**} \\
\\ 
\texttt{* في كل دور، يمكنك إرسال أمر واحد فقط من الأمرين التاليين:} \\
\texttt{ 1. `قل: <رسالة>`: لإرسال رسالة (كل شيء حتى فاصل السطر التالي) إلى اللاعب الآخر. سأقوم بتمريرها إلى شريكك.} \\
\texttt{2. `حرّك: <كائن>, (<X>, <Y>)`: لتحريك كائن إلى موضع جديد، حيث `<X>` هو العمود و `<Y>` هو الصف. سأخبرك بما إذا كانت حركتك ناجحة أم لا.} \\
\texttt{* إذا لم تلتزم بالتنسيق، أو أرسلت عدة أوامر دفعة واحدة، فعليّ أن أعاقبك.} \\
\texttt{* إذا تجاوز كلا اللاعبين أكثر من 12 عقوبة، فستخسران اللعبة معاً.} \\
\texttt{* من الضروري أن تتواصل مع اللاعب الآخر بشأن حالة هدفك! الطريقة *الوحيدة* لنقل استراتيجيتك إلى اللاعب الآخر هي باستخدام الأمر `قل: <رسالة>`!} \\
\\ 
\\ 
\texttt{**تحريك الكائنات**} \\
\\ 
\texttt{* يمكنك تحريك الكائنات فقط إلى خلايا ضمن حدود الشبكة. يجب أن تكون الخلية المستهدفة فارغة، أي يجب أن تحتوي فقط على الرمز '◌'.} \\
\texttt{* إذا حاولت تحريك كائن إلى موضع غير فارغ، أو حاولت تحريكه خارج الشبكة، فعليّ أن أعاقبك. ستحصل على محاولة أخرى.} \\
\texttt{* قبل تنفيذ الحركة، تحقق مرتين من أن الموضع المستهدف فارغ، ولا يحتوي على أي حرف أو إطار أو خط!} \\
\\ 
\\ 
\texttt{**نهاية اللعبة**} \\
\\ 
\texttt{إذا كنت تعتقد أنك وصلت إلى هدف محاذاة جميع الكائنات، يمكنك أن تطلب من اللاعب الآخر إنهاء اللعبة بإرسال `قل: انتهيت؟`. إذا طلب منك اللاعب الآخر إنهاء اللعبة، ورددت بـ `قل: انتهيت!`، ستنتهي اللعبة.} \\
\\ 
\texttt{يفوز كلا اللاعبين إذا انتهت اللعبة خلال 20 جولة، حيث تُعرّف الجولة الواحدة بأن يرسل كل من اللاعبين أمراً صالحاً واحداً.} \\
\\ 
\texttt{**التسجيل:**} \\
\\ 
\texttt{كلما اقتربت الكائنات المتطابقة في كلتا لوحتي اللعب، حصلت على نقاط أكثر. العقوبات تقلل نقاطك. هل يمكنك تحطيم الرقم القياسي؟} \\
\\ 
\\ 
\texttt{يرجى إرسال رسالة إلى اللاعب الآخر لبدء اللعبة!} \\
            }
        }
    }
    & \\ \\

    \theutterance \stepcounter{utterance}  
    & \multicolumn{4}{p{0.6\linewidth}}{
        \cellcolor[rgb]{0.9,0.9,0.9}{
            \makecell[{{p{\linewidth}}}]{
                \texttt{\tiny{[P1$\rangle$GM]}}
                \texttt{لنبدأ اللعبة بإرسال رسالة إلى اللاعب الآخر لنناقش استراتيجية مشتركة. سأبدأ باقتراح خطة لتحديد مواقع الكائنات المتطابقة وتحديد هدف مشترك.} \\
\\ 
\texttt{**رسالة أولى:**} \\
\texttt{`قل: مرحباً! لنبدأ بتحديد مواقع الكائنات المتطابقة. لديّ الكائنات التالية: W في (1,3)، I في (2,3)، T في (3,1)، H في (4,5)، C في (7,5). ما هي مواقع كائناتك؟ لنقارن ونقرر هدفاً مشتركاً.`} \\
\\ 
\texttt{سأرسل هذه الرسالة الآن، وبعد ذلك سننتظر رد اللاعب الآخر لنواصل التخطيط.} \\
            }
        }
    }
    & & \\ \\

    \theutterance \stepcounter{utterance}  
    & & & \multicolumn{2}{p{0.3\linewidth}}{
        \cellcolor[rgb]{0.9,0.9,0.9}{
            \makecell[{{p{\linewidth}}}]{
                \texttt{\tiny{[GM$|$GM]}}
                \texttt{لنبدأ اللعبة بإرسال رسالة إلى اللاعب الآخر لنناقش استراتيجية مشتركة. سأبدأ باقتراح خطة لتحديد مواقع الكائنات المتطابقة وتحديد هدف مشترك.} \\
\\ 
\texttt{**رسالة أولى:**} \\
\texttt{`قل: مرحباً! لنبدأ بتحديد مواقع الكائنات المتطابقة. لديّ الكائنات التالية: W في (1,3)، I في (2,3)، T في (3,1)، H في (4,5)، C في (7,5). ما هي مواقع كائناتك؟ لنقارن ونقرر هدفاً مشتركاً.`} \\
\\ 
\texttt{سأرسل هذه الرسالة الآن، وبعد ذلك سننتظر رد اللاعب الآخر لنواصل التخطيط.} \\
            }
        }
    }
    & & \\ \\

    \theutterance \stepcounter{utterance}  
    & & & \multicolumn{2}{p{0.3\linewidth}}{
        \cellcolor[rgb]{0.9,0.9,0.9}{
            \makecell[{{p{\linewidth}}}]{
                \texttt{\tiny{[GM$|$GM]}}
                \texttt{Invalid format: head and tail are not empty} \\
\texttt{head: 'لنبدأ اللعبة بإرسال رسالة إلى اللاعب الآخر لنناقش استراتيجية مشتركة. سأبدأ باقتراح خطة لتحديد مواقع الكائنات المتطابقة وتحديد هدف مشترك.} \\
\\ 
\texttt{**رسالة أولى:**} \\
\texttt{'} \\
\texttt{tail: '} \\
\\ 
\texttt{سأرسل هذه الرسالة الآن، وبعد ذلك سننتظر رد اللاعب الآخر لنواصل التخطيط.'} \\
            }
        }
    }
    & & \\ \\

    \theutterance \stepcounter{utterance}  
    & & \multicolumn{4}{p{0.6\linewidth}}{
        \cellcolor[rgb]{0.9,0.9,0.9}{
            \makecell[{{p{\linewidth}}}]{
                \texttt{\tiny{[P1$\langle$GM]}}
                \texttt{عقوبة: يجب ألا تحتوي رسالتك على أي شيء قبل الأمر أو بعده!} \\
\\ 
\texttt{تأكد من أن ردك يحتوي فقط على `قل: <رسالة>` أو `حرّك: <كائن>, (<X>, <Y>)`، ولا شيء غير ذلك!} \\
\\ 
\texttt{لقد جمعتما معاً 1 من أصل 12 عقوبة.} \\
\texttt{يرجى المحاولة مرة أخرى!} \\
            }
        }
    }
    & \\ \\

    \theutterance \stepcounter{utterance}  
    & \multicolumn{4}{p{0.6\linewidth}}{
        \cellcolor[rgb]{0.9,0.9,0.9}{
            \makecell[{{p{\linewidth}}}]{
                \texttt{\tiny{[P1$\rangle$GM]}}
                \texttt{`قل: لدي W في (1,3)، I في (2,3)، T في (3,1)، H في (4,5)، C في (7,5). ما هي مواقع كائناتك؟`} \\
            }
        }
    }
    & & \\ \\

    \theutterance \stepcounter{utterance}  
    & & & \multicolumn{2}{p{0.3\linewidth}}{
        \cellcolor[rgb]{0.9,0.9,0.9}{
            \makecell[{{p{\linewidth}}}]{
                \texttt{\tiny{[GM$|$GM]}}
                \texttt{`قل: لدي W في (1,3)، I في (2,3)، T في (3,1)، H في (4,5)، C في (7,5). ما هي مواقع كائناتك؟`} \\
            }
        }
    }
    & & \\ \\

    \theutterance \stepcounter{utterance}  
    & & \multicolumn{4}{p{0.6\linewidth}}{
        \cellcolor[rgb]{0.9,0.9,0.9}{
            \makecell[{{p{\linewidth}}}]{
                \texttt{\tiny{[GM$\rangle$P2]}}
                \texttt{أنا مدير لعبتك، وأنت تلعب لعبة تعاونية مع الشبكة التالية كلوحة لعب:} \\
\\ 
\texttt{```} \\
\texttt{ 1234567} \\
\texttt{╔══╤══╤═╗} \\
\texttt{║◌◌│W◌│I║ 1} \\
\texttt{╟──┼──┴─╢ 2} \\
\texttt{║◌◌│◌◌◌◌║ 3} \\
\texttt{╟──┘T◌◌◌║ 4} \\
\texttt{║H◌◌◌◌◌◌║ 5} \\
\texttt{║◌◌┌─┬──╢ 6} \\
\texttt{║C◌│◌│◌◌║ 7} \\
\texttt{╚══╧═╧══╝} \\
\\ 
\texttt{```} \\
\\ 
\texttt{* تعرض الحافة العلوية إحداثيات x التي تتزايد نحو اليمين، والحافة اليمنى إحداثيات y التي تتزايد نحو الأسفل.} \\
\texttt{* الكائنات التالية موضوعة عشوائياً على شبكتك: 'W', 'I', 'T', 'C', 'H'.} \\
\\ 
\\ 
\texttt{يرى اللاعب الآخر نسخة مختلفة من لوحة اللعب، حيث توضع الكائنات في مواقع عشوائية مختلفة. لا يمكنك رؤية لوحة اللاعب الآخر، ولا يمكنه رؤية لوحتك.} \\
\\ 
\\ 
\texttt{**الهدف:**} \\
\\ 
\texttt{يجب على كلا اللاعبين تحريك الكائنات على خلفية كل منهما بحيث تنتهي الكائنات المتطابقة عند الإحداثيات نفسها. عليك التواصل مع اللاعب الآخر للاتفاق على إعداد هدف مشترك.} \\
\\ 
\\ 
\texttt{**القواعد:**} \\
\\ 
\texttt{* في كل دور، يمكنك إرسال أمر واحد فقط من الأمرين التاليين:} \\
\texttt{ 1. `قل: <رسالة>`: لإرسال رسالة (كل شيء حتى فاصل السطر التالي) إلى اللاعب الآخر. سأقوم بتمريرها إلى شريكك.} \\
\texttt{2. `حرّك: <كائن>, (<X>, <Y>)`: لتحريك كائن إلى موضع جديد، حيث `<X>` هو العمود و `<Y>` هو الصف. سأخبرك بما إذا كانت حركتك ناجحة أم لا.} \\
\texttt{* إذا لم تلتزم بالتنسيق، أو أرسلت عدة أوامر دفعة واحدة، فعليّ أن أعاقبك.} \\
\texttt{* إذا تجاوز كلا اللاعبين أكثر من 12 عقوبة، فستخسران اللعبة معاً.} \\
\texttt{* من الضروري أن تتواصل مع اللاعب الآخر بشأن حالة هدفك! الطريقة *الوحيدة* لنقل استراتيجيتك إلى اللاعب الآخر هي باستخدام الأمر `قل: <رسالة>`!} \\
\\ 
\\ 
\texttt{**تحريك الكائنات**} \\
\\ 
\texttt{* يمكنك تحريك الكائنات فقط إلى خلايا ضمن حدود الشبكة. يجب أن تكون الخلية المستهدفة فارغة، أي يجب أن تحتوي فقط على الرمز '◌'.} \\
\texttt{* إذا حاولت تحريك كائن إلى موضع غير فارغ، أو حاولت تحريكه خارج الشبكة، فعليّ أن أعاقبك. ستحصل على محاولة أخرى.} \\
\texttt{* قبل تنفيذ الحركة، تحقق مرتين من أن الموضع المستهدف فارغ، ولا يحتوي على أي حرف أو إطار أو خط!} \\
\\ 
\\ 
\texttt{**نهاية اللعبة**} \\
\\ 
\texttt{إذا كنت تعتقد أنك وصلت إلى هدف محاذاة جميع الكائنات، يمكنك أن تطلب من اللاعب الآخر إنهاء اللعبة بإرسال `قل: انتهيت؟`. إذا طلب منك اللاعب الآخر إنهاء اللعبة، ورددت بـ `قل: انتهيت!`، ستنتهي اللعبة.} \\
\\ 
\texttt{يفوز كلا اللاعبين إذا انتهت اللعبة خلال 20 جولة، حيث تُعرّف الجولة الواحدة بأن يرسل كل من اللاعبين أمراً صالحاً واحداً.} \\
\\ 
\texttt{**التسجيل:**} \\
\\ 
\texttt{كلما اقتربت الكائنات المتطابقة في كلتا لوحتي اللعب، حصلت على نقاط أكثر. العقوبات تقلل نقاطك. هل يمكنك تحطيم الرقم القياسي؟} \\
\\ 
\texttt{بدأ اللاعب الآخر اللعبة بإرسال هذه الرسالة:} \\
\\ 
\texttt{"لدي W في (1,3)، I في (2,3)، T في (3,1)، H في (4,5)، C في (7,5). ما هي مواقع كائناتك؟"} \\
\\ 
\texttt{ما هو أمرك الأول؟} \\
            }
        }
    }
    & \\ \\

    \theutterance \stepcounter{utterance}  
    & & & \multicolumn{4}{p{0.6\linewidth}}{
        \cellcolor[rgb]{0.9,0.9,0.9}{
            \makecell[{{p{\linewidth}}}]{
                \texttt{\tiny{[GM$\langle$P2]}}
                \texttt{إليك اقتراحًا لأمرك الأول، والذي يجمع بين التواصل والحركة الاستراتيجية:} \\
\\ 
\texttt{**أمر:**} \\
\texttt{```} \\
\texttt{قل: لدي W في (2,1)، I في (6,1)، T في (3,4)، H في (1,5)، C في (1,7). لنبدأ بتحريك الكائنات المتطابقة نحو مواقع وسطية مشتركة. أقترح أن نضع W في (4,3)، I في (5,3)، T في (4,2)، H في (3,5)، C في (4,6). هل توافق؟} \\
\texttt{```} \\
\\ 
\texttt{\#\#\# سبب هذا الاقتراح:} \\
\texttt{1. **الشفافية**: تشارك مواقعك الحالية بوضوح لتجنب الالتباس.} \\
\texttt{2. **التنسيق**: تقترح مواقع وسطية مشتركة (مثل `(4,3)` لـ W) تجعل من السهل على كلا اللاعبين التحقق من المحاذاة دون الحاجة إلى تحركات معقدة.} \\
\texttt{3. **التحقق**: المواقع المقترحة فارغة في لوحتي اللعب (استنادًا إلى الشبكة المرئية)، مما يقلل من خطر العقوبات.} \\
\texttt{4. **الفعالية**: إذا وافق اللاعب الآخر، يمكنكما البدء في التحركات مباشرة في الجولات التالية.} \\
\\ 
\texttt{\#\#\# البديل (إذا كنت تفضل الحركة أولاً):} \\
\texttt{إذا كنت تريد تحريك كائن قبل الموافقة على الخطة، يمكنك البدء بـ:} \\
\texttt{```} \\
\texttt{حرّك: W, (4,3)} \\
\texttt{```} \\
\texttt{ثم إرسال الرسالة أعلاه في الجولة التالية. لكن التواصل المسبق يضمن توافقًا أفضل.} \\
\\ 
\texttt{**ملاحظة**: تأكد من أن المواقع المقترحة فارغة في لوحتي اللعب (مثل `(4,3)` و`(5,3)`) قبل اقتراحها. إذا كانت هناك تعارضات، اضبط الخطة بناءً على رد اللاعب الآخر.} \\
            }
        }
    }
     \\ \\

    \theutterance \stepcounter{utterance}  
    & & & \multicolumn{2}{p{0.3\linewidth}}{
        \cellcolor[rgb]{0.9,0.9,0.9}{
            \makecell[{{p{\linewidth}}}]{
                \texttt{\tiny{[GM$|$GM]}}
                \texttt{إليك اقتراحًا لأمرك الأول، والذي يجمع بين التواصل والحركة الاستراتيجية:} \\
\\ 
\texttt{**أمر:**} \\
\texttt{```} \\
\texttt{قل: لدي W في (2,1)، I في (6,1)، T في (3,4)، H في (1,5)، C في (1,7). لنبدأ بتحريك الكائنات المتطابقة نحو مواقع وسطية مشتركة. أقترح أن نضع W في (4,3)، I في (5,3)، T في (4,2)، H في (3,5)، C في (4,6). هل توافق؟} \\
\texttt{```} \\
\\ 
\texttt{\#\#\# سبب هذا الاقتراح:} \\
\texttt{1. **الشفافية**: تشارك مواقعك الحالية بوضوح لتجنب الالتباس.} \\
\texttt{2. **التنسيق**: تقترح مواقع وسطية مشتركة (مثل `(4,3)` لـ W) تجعل من السهل على كلا اللاعبين التحقق من المحاذاة دون الحاجة إلى تحركات معقدة.} \\
\texttt{3. **التحقق**: المواقع المقترحة فارغة في لوحتي اللعب (استنادًا إلى الشبكة المرئية)، مما يقلل من خطر العقوبات.} \\
\texttt{4. **الفعالية**: إذا وافق اللاعب الآخر، يمكنكما البدء في التحركات مباشرة في الجولات التالية.} \\
\\ 
\texttt{\#\#\# البديل (إذا كنت تفضل الحركة أولاً):} \\
\texttt{إذا كنت تريد تحريك كائن قبل الموافقة على الخطة، يمكنك البدء بـ:} \\
\texttt{```} \\
\texttt{حرّك: W, (4,3)} \\
\texttt{```} \\
\texttt{ثم إرسال الرسالة أعلاه في الجولة التالية. لكن التواصل المسبق يضمن توافقًا أفضل.} \\
\\ 
\texttt{**ملاحظة**: تأكد من أن المواقع المقترحة فارغة في لوحتي اللعب (مثل `(4,3)` و`(5,3)`) قبل اقتراحها. إذا كانت هناك تعارضات، اضبط الخطة بناءً على رد اللاعب الآخر.} \\
            }
        }
    }
    & & \\ \\

    \theutterance \stepcounter{utterance}  
    & & & \multicolumn{2}{p{0.3\linewidth}}{
        \cellcolor[rgb]{0.9,0.9,0.9}{
            \makecell[{{p{\linewidth}}}]{
                \texttt{\tiny{[GM$|$GM]}}
                \texttt{Invalid response format: إليك اقتراحًا لأمرك الأول، والذي يجمع بين التواصل والحركة الاستراتيجية:} \\
\\ 
\texttt{**أمر:**} \\
\\ 
\texttt{قل: لدي W في (2,1)، I في (6,1)، T في (3,4)، H في (1,5)، C في (1,7). لنبدأ بتحريك الكائنات المتطابقة نحو مواقع وسطية مشتركة. أقترح أن نضع W في (4,3)، I في (5,3)، T في (4,2)، H في (3,5)، C في (4,6). هل توافق؟} \\
\\ 
\\ 
\texttt{\#\#\# سبب هذا الاقتراح:} \\
\texttt{1. **الشفافية**: تشارك مواقعك الحالية بوضوح لتجنب الالتباس.} \\
\texttt{2. **التنسيق**: تقترح مواقع وسطية مشتركة (مثل (4,3) لـ W) تجعل من السهل على كلا اللاعبين التحقق من المحاذاة دون الحاجة إلى تحركات معقدة.} \\
\texttt{3. **التحقق**: المواقع المقترحة فارغة في لوحتي اللعب (استنادًا إلى الشبكة المرئية)، مما يقلل من خطر العقوبات.} \\
\texttt{4. **الفعالية**: إذا وافق اللاعب الآخر، يمكنكما البدء في التحركات مباشرة في الجولات التالية.} \\
\\ 
\texttt{\#\#\# البديل (إذا كنت تفضل الحركة أولاً):} \\
\texttt{إذا كنت تريد تحريك كائن قبل الموافقة على الخطة، يمكنك البدء بـ:} \\
\\ 
\texttt{حرّك: W, (4,3)} \\
\\ 
\texttt{ثم إرسال الرسالة أعلاه في الجولة التالية. لكن التواصل المسبق يضمن توافقًا أفضل.} \\
\\ 
\texttt{**ملاحظة**: تأكد من أن المواقع المقترحة فارغة في لوحتي اللعب (مثل (4,3) و(5,3)) قبل اقتراحها. إذا كانت هناك تعارضات، اضبط الخطة بناءً على رد اللاعب الآخر.} \\
            }
        }
    }
    & & \\ \\

    \theutterance \stepcounter{utterance}  
    & & \multicolumn{4}{p{0.6\linewidth}}{
        \cellcolor[rgb]{0.9,0.9,0.9}{
            \makecell[{{p{\linewidth}}}]{
                \texttt{\tiny{[GM$\rangle$P2]}}
                \texttt{عقوبة: تحتوي رسالتك على أكثر من أمر واحد!} \\
\\ 
\texttt{تأكد من أن ردك يحتوي فقط على `قل: <رسالة>` أو `حرّك: <كائن>, (<X>, <Y>)`، ولا شيء غير ذلك!} \\
\\ 
\texttt{لقد جمعتما معاً 2 من أصل 12 عقوبة.} \\
\texttt{يرجى المحاولة مرة أخرى!} \\
            }
        }
    }
    & \\ \\

    \theutterance \stepcounter{utterance}  
    & & & \multicolumn{4}{p{0.6\linewidth}}{
        \cellcolor[rgb]{0.9,0.9,0.9}{
            \makecell[{{p{\linewidth}}}]{
                \texttt{\tiny{[GM$\langle$P2]}}
                \texttt{شكرًا على التنبيه! سألتزم بالقواعد الآن.} \\
\\ 
\texttt{**أمر:**} \\
\texttt{```} \\
\texttt{قل: لدي W في (2,1)، I في (6,1)، T في (3,4)، H في (1,5)، C في (1,7). ما هي خطتك لتحريك الكائنات؟} \\
\texttt{```} \\
\\ 
\texttt{سأحتفظ باقتراحي للمواقع المشتركة في الجولة التالية بعد رد اللاعب الآخر.} \\
            }
        }
    }
     \\ \\

    \theutterance \stepcounter{utterance}  
    & & & \multicolumn{2}{p{0.3\linewidth}}{
        \cellcolor[rgb]{0.9,0.9,0.9}{
            \makecell[{{p{\linewidth}}}]{
                \texttt{\tiny{[GM$|$GM]}}
                \texttt{شكرًا على التنبيه! سألتزم بالقواعد الآن.} \\
\\ 
\texttt{**أمر:**} \\
\texttt{```} \\
\texttt{قل: لدي W في (2,1)، I في (6,1)، T في (3,4)، H في (1,5)، C في (1,7). ما هي خطتك لتحريك الكائنات؟} \\
\texttt{```} \\
\\ 
\texttt{سأحتفظ باقتراحي للمواقع المشتركة في الجولة التالية بعد رد اللاعب الآخر.} \\
            }
        }
    }
    & & \\ \\

    \theutterance \stepcounter{utterance}  
    & & & \multicolumn{2}{p{0.3\linewidth}}{
        \cellcolor[rgb]{0.9,0.9,0.9}{
            \makecell[{{p{\linewidth}}}]{
                \texttt{\tiny{[GM$|$GM]}}
                \texttt{Invalid format: head and tail are not empty} \\
\texttt{head: 'شكرًا على التنبيه! سألتزم بالقواعد الآن.} \\
\\ 
\texttt{**أمر:**} \\
\\ 
\texttt{'} \\
\texttt{tail: '} \\
\\ 
\\ 
\texttt{سأحتفظ باقتراحي للمواقع المشتركة في الجولة التالية بعد رد اللاعب الآخر.'} \\
            }
        }
    }
    & & \\ \\

    \theutterance \stepcounter{utterance}  
    & & \multicolumn{4}{p{0.6\linewidth}}{
        \cellcolor[rgb]{0.9,0.9,0.9}{
            \makecell[{{p{\linewidth}}}]{
                \texttt{\tiny{[GM$\rangle$P2]}}
                \texttt{عقوبة: يجب ألا تحتوي رسالتك على أي شيء قبل الأمر أو بعده!} \\
\\ 
\texttt{تأكد من أن ردك يحتوي فقط على `قل: <رسالة>` أو `حرّك: <كائن>, (<X>, <Y>)`، ولا شيء غير ذلك!} \\
\\ 
\texttt{لقد جمعتما معاً 3 من أصل 12 عقوبة.} \\
\texttt{يرجى المحاولة مرة أخرى!} \\
            }
        }
    }
    & \\ \\

    \theutterance \stepcounter{utterance}  
    & & & \multicolumn{4}{p{0.6\linewidth}}{
        \cellcolor[rgb]{0.9,0.9,0.9}{
            \makecell[{{p{\linewidth}}}]{
                \texttt{\tiny{[GM$\langle$P2]}}
                \texttt{أعتذر عن الخطأ. سألتزم بالتنسيق بدقة هذه المرة.} \\
\\ 
\texttt{**أمر:**} \\
\texttt{```} \\
\texttt{قل: مواقعي: W(2,1), I(6,1), T(3,4), H(1,5), C(1,7). أين كائناتك؟} \\
\texttt{```} \\
            }
        }
    }
     \\ \\

    \theutterance \stepcounter{utterance}  
    & & & \multicolumn{2}{p{0.3\linewidth}}{
        \cellcolor[rgb]{0.9,0.9,0.9}{
            \makecell[{{p{\linewidth}}}]{
                \texttt{\tiny{[GM$|$GM]}}
                \texttt{أعتذر عن الخطأ. سألتزم بالتنسيق بدقة هذه المرة.} \\
\\ 
\texttt{**أمر:**} \\
\texttt{```} \\
\texttt{قل: مواقعي: W(2,1), I(6,1), T(3,4), H(1,5), C(1,7). أين كائناتك؟} \\
\texttt{```} \\
            }
        }
    }
    & & \\ \\

    \theutterance \stepcounter{utterance}  
    & & & \multicolumn{2}{p{0.3\linewidth}}{
        \cellcolor[rgb]{0.9,0.9,0.9}{
            \makecell[{{p{\linewidth}}}]{
                \texttt{\tiny{[GM$|$GM]}}
                \texttt{Invalid format: head is not empty: 'أعتذر عن الخطأ. سألتزم بالتنسيق بدقة هذه المرة.} \\
\\ 
\texttt{**أمر:**} \\
\\ 
\texttt{'} \\
            }
        }
    }
    & & \\ \\

    \theutterance \stepcounter{utterance}  
    & & \multicolumn{4}{p{0.6\linewidth}}{
        \cellcolor[rgb]{0.9,0.9,0.9}{
            \makecell[{{p{\linewidth}}}]{
                \texttt{\tiny{[GM$\rangle$P2]}}
                \texttt{عقوبة: يجب ألا تحتوي رسالتك على أي شيء قبل الأمر!} \\
\\ 
\texttt{تأكد من أن ردك يحتوي فقط على `قل: <رسالة>` أو `حرّك: <كائن>, (<X>, <Y>)`، ولا شيء غير ذلك!} \\
\\ 
\texttt{لقد جمعتما معاً 4 من أصل 12 عقوبة.} \\
\texttt{يرجى المحاولة مرة أخرى!} \\
            }
        }
    }
    & \\ \\

    \theutterance \stepcounter{utterance}  
    & & & \multicolumn{4}{p{0.6\linewidth}}{
        \cellcolor[rgb]{0.9,0.9,0.9}{
            \makecell[{{p{\linewidth}}}]{
                \texttt{\tiny{[GM$\langle$P2]}}
                \texttt{فهمت! سأرسل الرسالة بشكل صحيح هذه المرة.} \\
\\ 
\texttt{**أمر:**} \\
\texttt{```} \\
\texttt{قل: W عند (2,1)، I عند (6,1)، T عند (3,4)، H عند (1,5)، C عند (1,7). أين كائناتك؟} \\
\texttt{```} \\
            }
        }
    }
     \\ \\

    \theutterance \stepcounter{utterance}  
    & & & \multicolumn{2}{p{0.3\linewidth}}{
        \cellcolor[rgb]{0.9,0.9,0.9}{
            \makecell[{{p{\linewidth}}}]{
                \texttt{\tiny{[GM$|$GM]}}
                \texttt{فهمت! سأرسل الرسالة بشكل صحيح هذه المرة.} \\
\\ 
\texttt{**أمر:**} \\
\texttt{```} \\
\texttt{قل: W عند (2,1)، I عند (6,1)، T عند (3,4)، H عند (1,5)، C عند (1,7). أين كائناتك؟} \\
\texttt{```} \\
            }
        }
    }
    & & \\ \\

    \theutterance \stepcounter{utterance}  
    & & & \multicolumn{2}{p{0.3\linewidth}}{
        \cellcolor[rgb]{0.9,0.9,0.9}{
            \makecell[{{p{\linewidth}}}]{
                \texttt{\tiny{[GM$|$GM]}}
                \texttt{Invalid format: head is not empty: 'فهمت! سأرسل الرسالة بشكل صحيح هذه المرة.} \\
\\ 
\texttt{**أمر:**} \\
\\ 
\texttt{'} \\
            }
        }
    }
    & & \\ \\

    \theutterance \stepcounter{utterance}  
    & & \multicolumn{4}{p{0.6\linewidth}}{
        \cellcolor[rgb]{0.9,0.9,0.9}{
            \makecell[{{p{\linewidth}}}]{
                \texttt{\tiny{[GM$\rangle$P2]}}
                \texttt{عقوبة: يجب ألا تحتوي رسالتك على أي شيء قبل الأمر!} \\
\\ 
\texttt{تأكد من أن ردك يحتوي فقط على `قل: <رسالة>` أو `حرّك: <كائن>, (<X>, <Y>)`، ولا شيء غير ذلك!} \\
\\ 
\texttt{لقد جمعتما معاً 5 من أصل 12 عقوبة.} \\
\texttt{يرجى المحاولة مرة أخرى!} \\
            }
        }
    }
    & \\ \\

    \theutterance \stepcounter{utterance}  
    & & & \multicolumn{4}{p{0.6\linewidth}}{
        \cellcolor[rgb]{0.9,0.9,0.9}{
            \makecell[{{p{\linewidth}}}]{
                \texttt{\tiny{[GM$\langle$P2]}}
                \texttt{**أمر:**} \\
\texttt{```} \\
\texttt{قل: مواقعي: W(2,1) I(6,1) T(3,4) H(1,5) C(1,7)} \\
\texttt{```} \\
            }
        }
    }
     \\ \\

    \theutterance \stepcounter{utterance}  
    & & & \multicolumn{2}{p{0.3\linewidth}}{
        \cellcolor[rgb]{0.9,0.9,0.9}{
            \makecell[{{p{\linewidth}}}]{
                \texttt{\tiny{[GM$|$GM]}}
                \texttt{**أمر:**} \\
\texttt{```} \\
\texttt{قل: مواقعي: W(2,1) I(6,1) T(3,4) H(1,5) C(1,7)} \\
\texttt{```} \\
            }
        }
    }
    & & \\ \\

    \theutterance \stepcounter{utterance}  
    & & & \multicolumn{2}{p{0.3\linewidth}}{
        \cellcolor[rgb]{0.9,0.9,0.9}{
            \makecell[{{p{\linewidth}}}]{
                \texttt{\tiny{[GM$|$GM]}}
                \texttt{Invalid format: head is not empty: '**أمر:**} \\
\\ 
\texttt{'} \\
            }
        }
    }
    & & \\ \\

    \theutterance \stepcounter{utterance}  
    & & \multicolumn{4}{p{0.6\linewidth}}{
        \cellcolor[rgb]{0.9,0.9,0.9}{
            \makecell[{{p{\linewidth}}}]{
                \texttt{\tiny{[GM$\rangle$P2]}}
                \texttt{عقوبة: يجب ألا تحتوي رسالتك على أي شيء قبل الأمر!} \\
\\ 
\texttt{تأكد من أن ردك يحتوي فقط على `قل: <رسالة>` أو `حرّك: <كائن>, (<X>, <Y>)`، ولا شيء غير ذلك!} \\
\\ 
\texttt{لقد جمعتما معاً 6 من أصل 12 عقوبة.} \\
\texttt{يرجى المحاولة مرة أخرى!} \\
            }
        }
    }
    & \\ \\

    \theutterance \stepcounter{utterance}  
    & & & \multicolumn{4}{p{0.6\linewidth}}{
        \cellcolor[rgb]{0.9,0.9,0.9}{
            \makecell[{{p{\linewidth}}}]{
                \texttt{\tiny{[GM$\langle$P2]}}
                \texttt{أعتذر عن الأخطاء المتكررة. سأرسل الرسالة بالتنسيق الصحيح الآن:} \\
\\ 
\texttt{**أمر:**} \\
\texttt{```} \\
\texttt{قل: W في (2,1)، I في (6,1)، T في (3,4)، H في (1,5)، C في (1,7)} \\
\texttt{```} \\
            }
        }
    }
     \\ \\

    \theutterance \stepcounter{utterance}  
    & & & \multicolumn{2}{p{0.3\linewidth}}{
        \cellcolor[rgb]{0.9,0.9,0.9}{
            \makecell[{{p{\linewidth}}}]{
                \texttt{\tiny{[GM$|$GM]}}
                \texttt{أعتذر عن الأخطاء المتكررة. سأرسل الرسالة بالتنسيق الصحيح الآن:} \\
\\ 
\texttt{**أمر:**} \\
\texttt{```} \\
\texttt{قل: W في (2,1)، I في (6,1)، T في (3,4)، H في (1,5)، C في (1,7)} \\
\texttt{```} \\
            }
        }
    }
    & & \\ \\

    \theutterance \stepcounter{utterance}  
    & & & \multicolumn{2}{p{0.3\linewidth}}{
        \cellcolor[rgb]{0.9,0.9,0.9}{
            \makecell[{{p{\linewidth}}}]{
                \texttt{\tiny{[GM$|$GM]}}
                \texttt{Invalid format: head is not empty: 'أعتذر عن الأخطاء المتكررة. سأرسل الرسالة بالتنسيق الصحيح الآن:} \\
\\ 
\texttt{**أمر:**} \\
\\ 
\texttt{'} \\
            }
        }
    }
    & & \\ \\

    \theutterance \stepcounter{utterance}  
    & & \multicolumn{4}{p{0.6\linewidth}}{
        \cellcolor[rgb]{0.9,0.9,0.9}{
            \makecell[{{p{\linewidth}}}]{
                \texttt{\tiny{[GM$\rangle$P2]}}
                \texttt{عقوبة: يجب ألا تحتوي رسالتك على أي شيء قبل الأمر!} \\
\\ 
\texttt{تأكد من أن ردك يحتوي فقط على `قل: <رسالة>` أو `حرّك: <كائن>, (<X>, <Y>)`، ولا شيء غير ذلك!} \\
\\ 
\texttt{لقد جمعتما معاً 7 من أصل 12 عقوبة.} \\
\texttt{يرجى المحاولة مرة أخرى!} \\
            }
        }
    }
    & \\ \\

    \theutterance \stepcounter{utterance}  
    & & & \multicolumn{4}{p{0.6\linewidth}}{
        \cellcolor[rgb]{0.9,0.9,0.9}{
            \makecell[{{p{\linewidth}}}]{
                \texttt{\tiny{[GM$\langle$P2]}}
                \texttt{فهمت الآن تمامًا {-} سأرسل الرسالة بدون أي كلمات تمهيدية:} \\
\\ 
\texttt{```} \\
\texttt{قل: W(2,1) I(6,1) T(3,4) H(1,5) C(1,7)} \\
\texttt{```} \\
            }
        }
    }
     \\ \\

    \theutterance \stepcounter{utterance}  
    & & & \multicolumn{2}{p{0.3\linewidth}}{
        \cellcolor[rgb]{0.9,0.9,0.9}{
            \makecell[{{p{\linewidth}}}]{
                \texttt{\tiny{[GM$|$GM]}}
                \texttt{فهمت الآن تمامًا {-} سأرسل الرسالة بدون أي كلمات تمهيدية:} \\
\\ 
\texttt{```} \\
\texttt{قل: W(2,1) I(6,1) T(3,4) H(1,5) C(1,7)} \\
\texttt{```} \\
            }
        }
    }
    & & \\ \\

    \theutterance \stepcounter{utterance}  
    & & & \multicolumn{2}{p{0.3\linewidth}}{
        \cellcolor[rgb]{0.9,0.9,0.9}{
            \makecell[{{p{\linewidth}}}]{
                \texttt{\tiny{[GM$|$GM]}}
                \texttt{Invalid format: head is not empty: 'فهمت الآن تمامًا {-} سأرسل الرسالة بدون أي كلمات تمهيدية:} \\
\\ 
\\ 
\texttt{'} \\
            }
        }
    }
    & & \\ \\

    \theutterance \stepcounter{utterance}  
    & & \multicolumn{4}{p{0.6\linewidth}}{
        \cellcolor[rgb]{0.9,0.9,0.9}{
            \makecell[{{p{\linewidth}}}]{
                \texttt{\tiny{[GM$\rangle$P2]}}
                \texttt{عقوبة: يجب ألا تحتوي رسالتك على أي شيء قبل الأمر!} \\
\\ 
\texttt{تأكد من أن ردك يحتوي فقط على `قل: <رسالة>` أو `حرّك: <كائن>, (<X>, <Y>)`، ولا شيء غير ذلك!} \\
\\ 
\texttt{لقد جمعتما معاً 8 من أصل 12 عقوبة.} \\
\texttt{يرجى المحاولة مرة أخرى!} \\
            }
        }
    }
    & \\ \\

    \theutterance \stepcounter{utterance}  
    & & & \multicolumn{4}{p{0.6\linewidth}}{
        \cellcolor[rgb]{0.9,0.9,0.9}{
            \makecell[{{p{\linewidth}}}]{
                \texttt{\tiny{[GM$\langle$P2]}}
                \texttt{أعتذر بشدة عن الإصرار على الخطأ. سأرسل الأمر بشكل صحيح تمامًا هذه المرة:} \\
\\ 
\texttt{```} \\
\texttt{W(2,1), I(6,1), T(3,4), H(1,5), C(1,7)} \\
\texttt{```} \\
\\ 
\texttt{**ملاحظة مهمة:** لقد نسيت كتابة `قل:` في البداية. الأمر الصحيح يجب أن يكون:} \\
\\ 
\texttt{```} \\
\texttt{قل: W(2,1), I(6,1), T(3,4), H(1,5), C(1,7)} \\
\texttt{```} \\
            }
        }
    }
     \\ \\

    \theutterance \stepcounter{utterance}  
    & & & \multicolumn{2}{p{0.3\linewidth}}{
        \cellcolor[rgb]{0.9,0.9,0.9}{
            \makecell[{{p{\linewidth}}}]{
                \texttt{\tiny{[GM$|$GM]}}
                \texttt{أعتذر بشدة عن الإصرار على الخطأ. سأرسل الأمر بشكل صحيح تمامًا هذه المرة:} \\
\\ 
\texttt{```} \\
\texttt{W(2,1), I(6,1), T(3,4), H(1,5), C(1,7)} \\
\texttt{```} \\
\\ 
\texttt{**ملاحظة مهمة:** لقد نسيت كتابة `قل:` في البداية. الأمر الصحيح يجب أن يكون:} \\
\\ 
\texttt{```} \\
\texttt{قل: W(2,1), I(6,1), T(3,4), H(1,5), C(1,7)} \\
\texttt{```} \\
            }
        }
    }
    & & \\ \\

    \theutterance \stepcounter{utterance}  
    & & & \multicolumn{2}{p{0.3\linewidth}}{
        \cellcolor[rgb]{0.9,0.9,0.9}{
            \makecell[{{p{\linewidth}}}]{
                \texttt{\tiny{[GM$|$GM]}}
                \texttt{Invalid format: head is not empty: 'أعتذر بشدة عن الإصرار على الخطأ. سأرسل الأمر بشكل صحيح تمامًا هذه المرة:} \\
\\ 
\\ 
\texttt{W(2,1), I(6,1), T(3,4), H(1,5), C(1,7)} \\
\\ 
\\ 
\texttt{**ملاحظة مهمة:** لقد نسيت كتابة قل: في البداية. الأمر الصحيح يجب أن يكون:} \\
\\ 
\\ 
\texttt{'} \\
            }
        }
    }
    & & \\ \\

    \theutterance \stepcounter{utterance}  
    & & \multicolumn{4}{p{0.6\linewidth}}{
        \cellcolor[rgb]{0.9,0.9,0.9}{
            \makecell[{{p{\linewidth}}}]{
                \texttt{\tiny{[GM$\rangle$P2]}}
                \texttt{عقوبة: يجب ألا تحتوي رسالتك على أي شيء قبل الأمر!} \\
\\ 
\texttt{تأكد من أن ردك يحتوي فقط على `قل: <رسالة>` أو `حرّك: <كائن>, (<X>, <Y>)`، ولا شيء غير ذلك!} \\
\\ 
\texttt{لقد جمعتما معاً 9 من أصل 12 عقوبة.} \\
\texttt{يرجى المحاولة مرة أخرى!} \\
            }
        }
    }
    & \\ \\

    \theutterance \stepcounter{utterance}  
    & & & \multicolumn{4}{p{0.6\linewidth}}{
        \cellcolor[rgb]{0.9,0.9,0.9}{
            \makecell[{{p{\linewidth}}}]{
                \texttt{\tiny{[GM$\langle$P2]}}
                \texttt{فهمت الآن تمامًا {-} سأرسل الأمر بالتنسيق الدقيق المطلوب دون أي إضافات:} \\
\\ 
\texttt{```} \\
\texttt{قل: مواقعي W(2,1) I(6,1) T(3,4) H(1,5) C(1,7)} \\
\texttt{```} \\
            }
        }
    }
     \\ \\

    \theutterance \stepcounter{utterance}  
    & & & \multicolumn{2}{p{0.3\linewidth}}{
        \cellcolor[rgb]{0.9,0.9,0.9}{
            \makecell[{{p{\linewidth}}}]{
                \texttt{\tiny{[GM$|$GM]}}
                \texttt{فهمت الآن تمامًا {-} سأرسل الأمر بالتنسيق الدقيق المطلوب دون أي إضافات:} \\
\\ 
\texttt{```} \\
\texttt{قل: مواقعي W(2,1) I(6,1) T(3,4) H(1,5) C(1,7)} \\
\texttt{```} \\
            }
        }
    }
    & & \\ \\

    \theutterance \stepcounter{utterance}  
    & & & \multicolumn{2}{p{0.3\linewidth}}{
        \cellcolor[rgb]{0.9,0.9,0.9}{
            \makecell[{{p{\linewidth}}}]{
                \texttt{\tiny{[GM$|$GM]}}
                \texttt{Invalid format: head is not empty: 'فهمت الآن تمامًا {-} سأرسل الأمر بالتنسيق الدقيق المطلوب دون أي إضافات:} \\
\\ 
\\ 
\texttt{'} \\
            }
        }
    }
    & & \\ \\

    \theutterance \stepcounter{utterance}  
    & & \multicolumn{4}{p{0.6\linewidth}}{
        \cellcolor[rgb]{0.9,0.9,0.9}{
            \makecell[{{p{\linewidth}}}]{
                \texttt{\tiny{[GM$\rangle$P2]}}
                \texttt{عقوبة: يجب ألا تحتوي رسالتك على أي شيء قبل الأمر!} \\
\\ 
\texttt{تأكد من أن ردك يحتوي فقط على `قل: <رسالة>` أو `حرّك: <كائن>, (<X>, <Y>)`، ولا شيء غير ذلك!} \\
\\ 
\texttt{لقد جمعتما معاً 10 من أصل 12 عقوبة.} \\
\texttt{يرجى المحاولة مرة أخرى!} \\
            }
        }
    }
    & \\ \\

    \theutterance \stepcounter{utterance}  
    & & & \multicolumn{4}{p{0.6\linewidth}}{
        \cellcolor[rgb]{0.9,0.9,0.9}{
            \makecell[{{p{\linewidth}}}]{
                \texttt{\tiny{[GM$\langle$P2]}}
                \texttt{أعتذر بشدة عن الأخطاء المتكررة. سأرسل الأمر النهائي بالتنسيق الصحيح تمامًا:} \\
\\ 
\texttt{```} \\
\texttt{قل: W في 2,1 I في 6,1 T في 3,4 H في 1,5 C في 1,7} \\
\texttt{```} \\
\\ 
\texttt{إذا كان هذا لا يزال غير صحيح، فالرجاء إخباري بالضبط بالتنسيق المطلوب دون أي أمثلة إضافية، وسأرسل الرسالة المطلوبة فقط دون أي إضافات.} \\
            }
        }
    }
     \\ \\

    \theutterance \stepcounter{utterance}  
    & & & \multicolumn{2}{p{0.3\linewidth}}{
        \cellcolor[rgb]{0.9,0.9,0.9}{
            \makecell[{{p{\linewidth}}}]{
                \texttt{\tiny{[GM$|$GM]}}
                \texttt{أعتذر بشدة عن الأخطاء المتكررة. سأرسل الأمر النهائي بالتنسيق الصحيح تمامًا:} \\
\\ 
\texttt{```} \\
\texttt{قل: W في 2,1 I في 6,1 T في 3,4 H في 1,5 C في 1,7} \\
\texttt{```} \\
\\ 
\texttt{إذا كان هذا لا يزال غير صحيح، فالرجاء إخباري بالضبط بالتنسيق المطلوب دون أي أمثلة إضافية، وسأرسل الرسالة المطلوبة فقط دون أي إضافات.} \\
            }
        }
    }
    & & \\ \\

    \theutterance \stepcounter{utterance}  
    & & & \multicolumn{2}{p{0.3\linewidth}}{
        \cellcolor[rgb]{0.9,0.9,0.9}{
            \makecell[{{p{\linewidth}}}]{
                \texttt{\tiny{[GM$|$GM]}}
                \texttt{Invalid format: head and tail are not empty} \\
\texttt{head: 'أعتذر بشدة عن الأخطاء المتكررة. سأرسل الأمر النهائي بالتنسيق الصحيح تمامًا:} \\
\\ 
\\ 
\texttt{'} \\
\texttt{tail: '} \\
\\ 
\\ 
\texttt{إذا كان هذا لا يزال غير صحيح، فالرجاء إخباري بالضبط بالتنسيق المطلوب دون أي أمثلة إضافية، وسأرسل الرسالة المطلوبة فقط دون أي إضافات.'} \\
            }
        }
    }
    & & \\ \\

    \theutterance \stepcounter{utterance}  
    & & \multicolumn{4}{p{0.6\linewidth}}{
        \cellcolor[rgb]{0.9,0.9,0.9}{
            \makecell[{{p{\linewidth}}}]{
                \texttt{\tiny{[GM$\rangle$P2]}}
                \texttt{عقوبة: يجب ألا تحتوي رسالتك على أي شيء قبل الأمر أو بعده!} \\
\\ 
\texttt{تأكد من أن ردك يحتوي فقط على `قل: <رسالة>` أو `حرّك: <كائن>, (<X>, <Y>)`، ولا شيء غير ذلك!} \\
\\ 
\texttt{لقد جمعتما معاً 11 من أصل 12 عقوبة.} \\
\texttt{يرجى المحاولة مرة أخرى!} \\
            }
        }
    }
    & \\ \\

    \theutterance \stepcounter{utterance}  
    & & & \multicolumn{4}{p{0.6\linewidth}}{
        \cellcolor[rgb]{0.9,0.9,0.9}{
            \makecell[{{p{\linewidth}}}]{
                \texttt{\tiny{[GM$\langle$P2]}}
                \texttt{```} \\
\texttt{قل: W(2,1) I(6,1) T(3,4) H(1,5) C(1,7)} \\
\texttt{```} \\
            }
        }
    }
     \\ \\

    \theutterance \stepcounter{utterance}  
    & & & \multicolumn{2}{p{0.3\linewidth}}{
        \cellcolor[rgb]{0.9,0.9,0.9}{
            \makecell[{{p{\linewidth}}}]{
                \texttt{\tiny{[GM$|$GM]}}
                \texttt{```} \\
\texttt{قل: W(2,1) I(6,1) T(3,4) H(1,5) C(1,7)} \\
\texttt{```} \\
            }
        }
    }
    & & \\ \\

    \theutterance \stepcounter{utterance}  
    & & \multicolumn{4}{p{0.6\linewidth}}{
        \cellcolor[rgb]{0.9,0.9,0.9}{
            \makecell[{{p{\linewidth}}}]{
                \texttt{\tiny{[P1$\langle$GM]}}
                \texttt{تم تمرير رسالتك إلى اللاعب الآخر.} \\
\\ 
\\ 
\texttt{أنت تلعب حالياً الجولة 1 من أصل 20 كحد أقصى.} \\
\texttt{لقد جمعتما معاً 11 من أصل 12 عقوبة.} \\
\texttt{أرسل اللاعب الآخر هذه الرسالة:} \\
\\ 
\\ 
\texttt{"W(2,1) I(6,1) T(3,4) H(1,5) C(1,7)"} \\
\\ 
\\ 
\texttt{ما هو أمرك التالي؟} \\
            }
        }
    }
    & \\ \\

    \theutterance \stepcounter{utterance}  
    & \multicolumn{4}{p{0.6\linewidth}}{
        \cellcolor[rgb]{0.9,0.9,0.9}{
            \makecell[{{p{\linewidth}}}]{
                \texttt{\tiny{[P1$\rangle$GM]}}
                \texttt{`قل: لنضع هدفاً مشتركاً: W→(2,1)، I→(6,1)، T→(3,4)، H→(1,5)، C→(1,7). سأبدأ بتحريك I إلى (6,3) مؤقتاً لتجنب التعارض. ماذا عنك؟`} \\
            }
        }
    }
    & & \\ \\

    \theutterance \stepcounter{utterance}  
    & & & \multicolumn{2}{p{0.3\linewidth}}{
        \cellcolor[rgb]{0.9,0.9,0.9}{
            \makecell[{{p{\linewidth}}}]{
                \texttt{\tiny{[GM$|$GM]}}
                \texttt{`قل: لنضع هدفاً مشتركاً: W→(2,1)، I→(6,1)، T→(3,4)، H→(1,5)، C→(1,7). سأبدأ بتحريك I إلى (6,3) مؤقتاً لتجنب التعارض. ماذا عنك؟`} \\
            }
        }
    }
    & & \\ \\

    \theutterance \stepcounter{utterance}  
    & & \multicolumn{4}{p{0.6\linewidth}}{
        \cellcolor[rgb]{0.9,0.9,0.9}{
            \makecell[{{p{\linewidth}}}]{
                \texttt{\tiny{[GM$\rangle$P2]}}
                \texttt{تم تمرير رسالتك إلى اللاعب الآخر.} \\
\\ 
\\ 
\texttt{أنت تلعب حالياً الجولة 2 من أصل 20 كحد أقصى.} \\
\texttt{لقد جمعتما معاً 11 من أصل 12 عقوبة.} \\
\texttt{أرسل اللاعب الآخر هذه الرسالة:} \\
\\ 
\\ 
\texttt{"لنضع هدفاً مشتركاً: W→(2,1)، I→(6,1)، T→(3,4)، H→(1,5)، C→(1,7). سأبدأ بتحريك I إلى (6,3) مؤقتاً لتجنب التعارض. ماذا عنك؟"} \\
\\ 
\\ 
\texttt{ما هو أمرك التالي؟} \\
            }
        }
    }
    & \\ \\

    \theutterance \stepcounter{utterance}  
    & & & \multicolumn{4}{p{0.6\linewidth}}{
        \cellcolor[rgb]{0.9,0.9,0.9}{
            \makecell[{{p{\linewidth}}}]{
                \texttt{\tiny{[GM$\langle$P2]}}
                \texttt{بما أننا اقتربنا من الحد الأقصى للعقوبات، سأكون حذرًا للغاية في الرد. سأختار تحريك كائن واحد فقط لتجنب أي أخطاء:} \\
\\ 
\texttt{```} \\
\texttt{حرّك: I, (6,3)} \\
\texttt{```} \\
\\ 
\texttt{هذا يتوافق مع اقتراح اللاعب الآخر بتحريك I مؤقتًا لتجنب التعارض. بعد تنفيذ هذه الحركة بنجاح، يمكننا التنسيق بشأن الخطوات التالية.} \\
            }
        }
    }
     \\ \\

    \theutterance \stepcounter{utterance}  
    & & & \multicolumn{2}{p{0.3\linewidth}}{
        \cellcolor[rgb]{0.9,0.9,0.9}{
            \makecell[{{p{\linewidth}}}]{
                \texttt{\tiny{[GM$|$GM]}}
                \texttt{بما أننا اقتربنا من الحد الأقصى للعقوبات، سأكون حذرًا للغاية في الرد. سأختار تحريك كائن واحد فقط لتجنب أي أخطاء:} \\
\\ 
\texttt{```} \\
\texttt{حرّك: I, (6,3)} \\
\texttt{```} \\
\\ 
\texttt{هذا يتوافق مع اقتراح اللاعب الآخر بتحريك I مؤقتًا لتجنب التعارض. بعد تنفيذ هذه الحركة بنجاح، يمكننا التنسيق بشأن الخطوات التالية.} \\
            }
        }
    }
    & & \\ \\

    \theutterance \stepcounter{utterance}  
    & & & \multicolumn{2}{p{0.3\linewidth}}{
        \cellcolor[rgb]{0.9,0.9,0.9}{
            \makecell[{{p{\linewidth}}}]{
                \texttt{\tiny{[GM$|$GM]}}
                \texttt{Invalid format: head and tail are not empty} \\
\texttt{head: 'بما أننا اقتربنا من الحد الأقصى للعقوبات، سأكون حذرًا للغاية في الرد. سأختار تحريك كائن واحد فقط لتجنب أي أخطاء:} \\
\\ 
\\ 
\texttt{'} \\
\texttt{tail: '} \\
\\ 
\\ 
\texttt{هذا يتوافق مع اقتراح اللاعب الآخر بتحريك I مؤقتًا لتجنب التعارض. بعد تنفيذ هذه الحركة بنجاح، يمكننا التنسيق بشأن الخطوات التالية.'} \\
            }
        }
    }
    & & \\ \\

    \theutterance \stepcounter{utterance}  
    & & \multicolumn{4}{p{0.6\linewidth}}{
        \cellcolor[rgb]{0.9,0.9,0.9}{
            \makecell[{{p{\linewidth}}}]{
                \texttt{\tiny{[GM$\rangle$P2]}}
                \texttt{عقوبة: يجب ألا تحتوي رسالتك على أي شيء قبل الأمر أو بعده!} \\
\\ 
\texttt{تأكد من أن ردك يحتوي فقط على `قل: <رسالة>` أو `حرّك: <كائن>, (<X>, <Y>)`، ولا شيء غير ذلك!} \\
\\ 
\texttt{لقد جمعتما معاً 12 من أصل 12 عقوبة.} \\
\texttt{يرجى المحاولة مرة أخرى!} \\
            }
        }
    }
    & \\ \\

    \theutterance \stepcounter{utterance}  
    & & & \multicolumn{4}{p{0.6\linewidth}}{
        \cellcolor[rgb]{0.9,0.9,0.9}{
            \makecell[{{p{\linewidth}}}]{
                \texttt{\tiny{[GM$\langle$P2]}}
                \texttt{```} \\
\texttt{حرّك: I,(6,3)} \\
\texttt{```} \\
\\ 
\texttt{(ملاحظة: أزلت المسافة بعد الفاصلة لتجنب أي مشاكل في التنسيق)} \\
            }
        }
    }
     \\ \\

    \theutterance \stepcounter{utterance}  
    & & & \multicolumn{2}{p{0.3\linewidth}}{
        \cellcolor[rgb]{0.9,0.9,0.9}{
            \makecell[{{p{\linewidth}}}]{
                \texttt{\tiny{[GM$|$GM]}}
                \texttt{```} \\
\texttt{حرّك: I,(6,3)} \\
\texttt{```} \\
\\ 
\texttt{(ملاحظة: أزلت المسافة بعد الفاصلة لتجنب أي مشاكل في التنسيق)} \\
            }
        }
    }
    & & \\ \\

    \theutterance \stepcounter{utterance}  
    & & & \multicolumn{2}{p{0.3\linewidth}}{
        \cellcolor[rgb]{0.9,0.9,0.9}{
            \makecell[{{p{\linewidth}}}]{
                \texttt{\tiny{[GM$|$GM]}}
                \texttt{Invalid response format} \\
            }
        }
    }
    & & \\ \\

    \theutterance \stepcounter{utterance}  
    & & & \multicolumn{2}{p{0.3\linewidth}}{
        \cellcolor[rgb]{0.9,0.9,0.9}{
            \makecell[{{p{\linewidth}}}]{
                \texttt{\tiny{[GM$|$GM]}}
                \texttt{Maximum number of penalties exceeded} \\
            }
        }
    }
    & & \\ \\

    \theutterance \stepcounter{utterance}  
    & & & \multicolumn{2}{p{0.3\linewidth}}{
        \cellcolor[rgb]{0.9,0.9,0.9}{
            \makecell[{{p{\linewidth}}}]{
                \texttt{\tiny{[GM$|$GM]}}
                \texttt{Initial states:} \\
\texttt{```} \\
\texttt{Player 1:} \\
\texttt{ 1234567} \\
\texttt{╔══╤══╤═╗} \\
\texttt{║◌◌│T◌│◌║ 1} \\
\texttt{╟──┼──┴─╢ 2} \\
\texttt{║WI│◌◌◌◌║ 3} \\
\texttt{╟──┘◌◌◌◌║ 4} \\
\texttt{║◌◌◌H◌◌C║ 5} \\
\texttt{║◌◌┌─┬──╢ 6} \\
\texttt{║◌◌│◌│◌◌║ 7} \\
\texttt{╚══╧═╧══╝} \\
\\ 
\\ 
\texttt{Player 2:} \\
\texttt{ 1234567} \\
\texttt{╔══╤══╤═╗} \\
\texttt{║◌◌│W◌│I║ 1} \\
\texttt{╟──┼──┴─╢ 2} \\
\texttt{║◌◌│◌◌◌◌║ 3} \\
\texttt{╟──┘T◌◌◌║ 4} \\
\texttt{║H◌◌◌◌◌◌║ 5} \\
\texttt{║◌◌┌─┬──╢ 6} \\
\texttt{║C◌│◌│◌◌║ 7} \\
\texttt{╚══╧═╧══╝} \\
\\ 
\texttt{```} \\
            }
        }
    }
    & & \\ \\

    \theutterance \stepcounter{utterance}  
    & & & \multicolumn{2}{p{0.3\linewidth}}{
        \cellcolor[rgb]{0.9,0.9,0.9}{
            \makecell[{{p{\linewidth}}}]{
                \texttt{\tiny{[GM$|$GM]}}
                \texttt{End states:} \\
\texttt{```} \\
\texttt{Player 1:} \\
\texttt{ 1234567} \\
\texttt{╔══╤══╤═╗} \\
\texttt{║◌◌│T◌│◌║ 1} \\
\texttt{╟──┼──┴─╢ 2} \\
\texttt{║WI│◌◌◌◌║ 3} \\
\texttt{╟──┘◌◌◌◌║ 4} \\
\texttt{║◌◌◌H◌◌C║ 5} \\
\texttt{║◌◌┌─┬──╢ 6} \\
\texttt{║◌◌│◌│◌◌║ 7} \\
\texttt{╚══╧═╧══╝} \\
\\ 
\\ 
\texttt{Player 2:} \\
\texttt{ 1234567} \\
\texttt{╔══╤══╤═╗} \\
\texttt{║◌◌│W◌│I║ 1} \\
\texttt{╟──┼──┴─╢ 2} \\
\texttt{║◌◌│◌◌◌◌║ 3} \\
\texttt{╟──┘T◌◌◌║ 4} \\
\texttt{║H◌◌◌◌◌◌║ 5} \\
\texttt{║◌◌┌─┬──╢ 6} \\
\texttt{║C◌│◌│◌◌║ 7} \\
\texttt{╚══╧═╧══╝} \\
\\ 
\texttt{```} \\
            }
        }
    }
    & & \\ \\

    \theutterance \stepcounter{utterance}  
    & & & \multicolumn{2}{p{0.3\linewidth}}{
        \cellcolor[rgb]{0.9,0.9,0.9}{
            \makecell[{{p{\linewidth}}}]{
                \texttt{\tiny{[GM$|$GM]}}
                \texttt{* success: False} \\
\texttt{* lose: False} \\
\texttt{* aborted: True} \\
\texttt{{-}{-}{-}{-}{-}{-}{-}} \\
\texttt{* Shifts: 0.00} \\
\texttt{* Max Shifts: 8.00} \\
\texttt{* Min Shifts: 4.00} \\
\texttt{* End Distance Sum: 21.32} \\
\texttt{* Init Distance Sum: 21.32} \\
\texttt{* Expected Distance Sum: 20.95} \\
\texttt{* Penalties: 13.00} \\
\texttt{* Max Penalties: 12.00} \\
\texttt{* Rounds: 1.00} \\
\texttt{* Max Rounds: 20.00} \\
\texttt{* Object Count: 5.00} \\
            }
        }
    }
    & & \\ \\

\end{supertabular}
}

\end{document}
